\section{Discussion}
\label{sec:discussion}

\subsection{The Convergence--Differentiation Tension}
\label{sec:tension}

Our most striking finding is that focal point training---which teaches agents to converge on the same answer---can \emph{destroy} coordination on tasks requiring complementary behavior.
The \focalpt game drops role assignment to 0\% on moving, resource, and timing tasks, because both agents learn to pick the ``obvious'' role.
This reveals a fundamental axis of coordination that any training intervention must address: some tasks require matching (focal points), while others require dividing (roles).

The \aidesigned game resolves this tension by explicitly training both skills.
Its system prompt includes not only focal point identification but also perspective-taking and role differentiation, teaching agents meta-cognitive awareness of \emph{which} coordination strategy to apply.
The result is 100\% complementarity on all hard role tasks while maintaining competitive focal point performance.

This finding has implications beyond our experimental setup.
Real-world coordination involves both convergence (\eg agreeing on meeting times, choosing compatible tools) and differentiation (\eg dividing labor, assigning complementary roles).
Training programs that focus on only one skill may inadvertently harm the other.

\subsection{Why Multi-Skill Training Works}
\label{sec:multiskill}

The \syncminds game combines four coordination skills:
(1) focal point identification---recognizing culturally salient choices;
(2) perspective-taking---predicting what a partner will do;
(3) role differentiation---choosing complementary rather than identical actions;
(4) convention formation---building shared protocols through interaction.

We hypothesize that this combination works because it provides agents with a richer \emph{coordination vocabulary}.
Single-skill games give agents one tool (``converge'' or ``differentiate''); the multi-skill game gives them a toolkit and, crucially, the meta-skill of selecting the right tool for the task.

\subsection{LLMs as Coordination Proxies}
\label{sec:proxies}

Our use of LLM agents as human proxies enables experiments at a scale (12,000 API calls, 50 pairs $\times$ 31 tasks $\times$ 5 conditions) that would be prohibitively expensive with human participants.
The approach is motivated by \citet{riedl2026emergent}'s finding that LLM agents exhibit emergent coordination similar to human coordination.

However, LLMs and humans coordinate through different mechanisms.
LLM coordination is driven by shared training data---cultural knowledge encoded in internet text---while human coordination additionally involves spatial reasoning, emotional salience, and embodied experience.
The near-ceiling baseline on standard focal point tasks (92.8\%) illustrates this: LLMs have extremely strong cultural priors that make simple coordination trivial.

We view LLM experiments as a \emph{first stage} that identifies promising game designs for subsequent human validation, not as a replacement for human studies.

\subsection{Limitations}
\label{sec:limitations}

\para{In-context learning vs.\ actual gameplay.}
We simulate game experience through system prompt descriptions rather than actual interactive gameplay.
Real games involve procedural learning, emotional engagement, repeated practice, and embodied experience, all of which may produce different (likely stronger) coordination effects.
Our results should be interpreted as a lower bound on the potential of game-based coordination training.

\para{Temperature sensitivity.}
All experiments use temperature $= 0.9$.
Lower temperatures would increase coordination by reducing response diversity (trivially), while higher temperatures might reveal different patterns.
A systematic temperature sweep would strengthen our conclusions.

\para{Single model.}
We test only \gptfour.
Different LLMs have different cultural knowledge distributions and may respond differently to game training prompts.
Cross-model replication (\eg with Claude, Gemini, or open-source models) is needed to establish generality.

\para{Normalization artifacts.}
Response matching depends on text normalization.
Some matches or mismatches may be artifacts of how we process responses (\eg whether ``NY'' and ``New York'' are treated as equivalent).

\para{Sample size.}
$N{=}50$ pairs per condition provides moderate statistical power.
Extreme results (0\% or 100\%) should be interpreted with appropriate caution, as they may reflect deterministic behavior at this temperature rather than robust population-level effects.

\para{Control condition variability.}
The \trivia control sometimes outperforms coordination-specific games on individual tasks (\eg year: 100\% vs.\ baseline 56\%), suggesting that any system prompt can shift LLM coordination patterns.
This makes it important to interpret treatment effects relative to both the no-game baseline and the active control.
